\documentclass[11pt]{article}

% Page geometry
\usepackage[margin=1in]{geometry}

% Standard packages
\usepackage[utf8]{inputenc}
\usepackage[T1]{fontenc}
\usepackage{amsmath,amssymb,amsfonts}
\usepackage{booktabs}
\usepackage{graphicx}
\usepackage{hyperref}
\usepackage{xcolor}
\usepackage{tabularx}
\usepackage{natbib}
% \usepackage{microtype} % Uncomment if full TeX Live is installed
\usepackage{url}
\usepackage{float}

\hypersetup{
  colorlinks=true,
  linkcolor=blue!70!black,
  citecolor=green!50!black,
  urlcolor=blue!60!black
}

% Compact spacing
\setlength{\parskip}{0.4em}
\setlength{\parindent}{0em}

\title{\textbf{OpenProblemAtlas}: A Machine-Verifiable Atlas of \\Unsolved Problems for AI Mathematical Reasoning}

\author{
  OpenProblemAtlas Contributors\thanks{Project website: \url{https://openproblem-atlas.org}. Data and code available at \url{https://github.com/openproblem-atlas/OpenProblemAtlas}.}
}

\date{}

\begin{document}

\maketitle

% ============================================================
\begin{abstract}
We present \textsc{OpenProblemAtlas} (OPA), the first large-scale, structured, and machine-verifiable dataset of \emph{unsolved} mathematical problems, designed to evaluate and drive AI systems toward genuine mathematical discovery rather than mere problem-solving on known answers.
OPA contains 2{,}075 open problems across three domains---mathematics, theoretical computer science, and mathematical physics---each annotated with a YAML-based data model that records problem status, domain taxonomy, a three-axis verification profile (statement precision, solution checkability, machine actionability), a tier classification, and five numerical scores measuring research fitness.
The dataset is accompanied by 2{,}075 verification contracts, 101 executable Python checkers, and 312 Lean\,4 formalization stubs covering 15\% of problems.
We introduce \textsc{OPA-Bench}, a curated benchmark of 100 problems graded at Bronze, Silver, and Gold difficulty levels for evaluating AI mathematical reasoning on open questions.
Unlike existing benchmarks such as MATH, GSM8K, and MiniF2F, where ground-truth answers are known, OPA targets the frontier of mathematical knowledge, providing infrastructure for machine-verifiable partial progress on unsolved problems.
All data, contracts, checkers, and formalizations are released under an open license.
\end{abstract}

% ============================================================
\section{Introduction}
\label{sec:intro}

Recent advances in AI-driven mathematical reasoning have produced systems capable of impressive performance on established benchmarks.
AlphaProof~\citep{alphaproof2024} achieved silver-medal-level performance at the International Mathematical Olympiad, while DeepSeek-Prover~\citep{xin2024deepseekprover} and Llemma~\citep{azerbayev2024llemma} have demonstrated strong theorem-proving capabilities in Lean and other formal systems.
AlphaGeometry~\citep{trinh2024alphageometry} solved olympiad geometry problems at a near-gold-medal level, and language-model-based approaches~\citep{polu2020gpt,jiang2023draft} have shown that large models can generate and guide formal proofs.

Yet all major benchmarks used to evaluate these systems---MATH~\citep{hendrycks2021math}, GSM8K~\citep{cobbe2021gsm8k}, MiniF2F~\citep{zheng2022minif2f}, PutnamBench~\citep{tsoukalas2024putnambench}---share a critical limitation: \textbf{the answers are already known}.
These benchmarks measure the ability of AI systems to reproduce solutions that humans have already found, not to discover new mathematical truths.
An AI system that scores perfectly on MATH has demonstrated sophisticated pattern matching and reasoning, but it has not advanced the frontier of mathematical knowledge.

This gap motivates \textsc{OpenProblemAtlas} (OPA).
We build the first structured, machine-readable dataset of \emph{unsolved} problems, where no answer is known a priori.
OPA is not merely a list of open problems (cf.\ Open Problem Garden~\citep{openproblemgarden}, Erd\H{o}s Problems~\citep{erdosproblems}); it is a \emph{verification infrastructure} that enables AI systems to make machine-checkable partial progress.
Each problem comes with:
\begin{enumerate}  \item A structured YAML entry with canonical statement, domain taxonomy, and status metadata;
  \item A \textbf{verification contract} specifying how claimed progress can be checked;
  \item Where applicable, an executable \textbf{Python checker} for computational verification;
  \item Where applicable, a \textbf{Lean\,4 formalization stub} for formal proof attempts.
\end{enumerate}

Our contributions are:
\begin{itemize}  \item A curated dataset of 2{,}075 open problems with rich structured annotations spanning three major domains;
  \item A multi-layer verification framework with 2{,}075 contracts and 101 executable checkers;
  \item 312 Lean\,4 formalization stubs (15\% coverage) integrated with the verification pipeline;
  \item \textsc{OPA-Bench}, a benchmark of 100 problems at three difficulty levels for standardized AI evaluation;
  \item An open-source, continuously updated data infrastructure with schema validation and CI/CD.
\end{itemize}

% ============================================================
\section{Related Work}
\label{sec:related}

\paragraph{Mathematical Reasoning Benchmarks.}
MATH~\citep{hendrycks2021math} provides 12{,}500 competition-level problems with final numerical or symbolic answers.
GSM8K~\citep{cobbe2021gsm8k} contains 8{,}500 grade-school word problems.
MiniF2F~\citep{zheng2022minif2f} offers 488 formalized olympiad-level statements in Lean, Isabelle, and Metamath.
PutnamBench~\citep{tsoukalas2024putnambench} contributes 1{,}697 Putnam competition problems formalized in Lean\,4.
All of these benchmarks evaluate systems on problems with \emph{known} solutions.

\paragraph{Open Problem Collections.}
Several informal collections of unsolved problems exist.
The Open Problem Garden~\citep{openproblemgarden} hosts a wiki-style database organized by topic.
Erd\H{o}s Problems~\citep{erdosproblems} catalogs problems posed by Paul Erd\H{o}s with their current status.
The Clay Millennium Problems represent seven high-profile open problems with \$1M prizes.
However, none of these resources provides machine-readable structured data, verification contracts, or executable checkers.

\paragraph{Formal Mathematics and Autoformalization.}
The Lean\,4 proof assistant~\citep{moura2021lean4} and its mathematical library mathlib4~\citep{mathlib2020} provide a growing corpus of formalized mathematics.
LeanDojo~\citep{yang2024leandojo} and related tools enable neural theorem provers to interact with Lean.
Recent work on autoformalization~\citep{wu2022autoformalization} and FormalConjectures~\citep{glazer2024formalconjectures} explores the automatic translation of informal mathematics into formal statements, which is complementary to OPA's formalization stubs.

% ============================================================
\section{Dataset Description}
\label{sec:dataset}

\subsection{Data Model}

Each problem in OPA is represented as a YAML file conforming to a JSON Schema (version 1.0.0).
The schema enforces the following required fields:

\begin{itemize}  \item \texttt{id}: A globally unique identifier following the pattern \texttt{opa.<domain>.<subdomain>.<slug>};
  \item \texttt{title}: Canonical problem name;
  \item \texttt{status}: A structured object with \texttt{label} (one of \texttt{open}, \texttt{partially\_solved}, \texttt{solved}, \texttt{disproved}, \texttt{conditional}, \texttt{independent}, \texttt{ambiguous}, \texttt{retired\_duplicate}), \texttt{confidence} level, review date, and review state;
  \item \texttt{domain} and \texttt{subdomains}: Hierarchical classification;
  \item \texttt{statement}: Canonical and informal descriptions;
  \item \texttt{verification\_profile}: Three-axis verification classification;
  \item \texttt{tier}: Tier assignment (\texttt{tier\_1}, \texttt{tier\_2}, or \texttt{tier\_3});
  \item \texttt{sources}: Bibliographic references with DOIs and URLs.
\end{itemize}

Optional fields include \texttt{problem\_type}, \texttt{answer\_type}, \texttt{partial\_results}, \texttt{relations} (equivalences, generalizations, special cases), \texttt{formalization} metadata, five-dimensional \texttt{scores}, and \texttt{prize} information.

\subsection{Verification Profile}
\label{sec:verif-profile}

A distinctive feature of OPA is the \textbf{three-axis verification profile}, which classifies each problem along three orthogonal dimensions:

\begin{enumerate}  \item \textbf{Statement Precision} (\texttt{high}/\texttt{medium}/\texttt{low}): How precisely the problem can be stated as a formal mathematical proposition. \emph{High}: a precise conjecture, bound, or existence claim. \emph{Medium}: clear mathematical objects but some definitional ambiguity. \emph{Low}: a research direction or vague agenda.
  \item \textbf{Solution Checkability} (\texttt{proof\_assistant}/\texttt{computational}/\texttt{mixed}/\texttt{expert\_review}): How a claimed solution or refutation can be verified. Problems checkable by proof assistants (Lean, Coq) or by computational means (SAT/SMT solvers, counterexample search) are the most machine-tractable.
  \item \textbf{Machine Actionability} (\texttt{high}/\texttt{medium}/\texttt{low}): How much AI agents can currently participate in solving the problem. \emph{High}: AI can directly attempt the problem via search, symbolic computation, or formalization. \emph{Low}: the problem currently lies beyond machine capabilities.
\end{enumerate}

\subsection{Tier System}

The tier system provides a coarse classification of machine tractability:

\begin{itemize}  \item \textbf{Tier 1} (175 problems): Core machine-actionable problems with high statement precision \emph{and} machine-checkable solutions (proof assistant or computational) \emph{and} at least medium machine actionability.
  \item \textbf{Tier 2} (106 problems): Extended problems with precise statements but limited machine verification---typically requiring mixed human--machine approaches.
  \item \textbf{Tier 3} (1{,}794 problems): Candidate problems that require further structuring; typically have low statement precision or require expert review for solution checking.
\end{itemize}

\subsection{Scoring Dimensions}

For a subset of 312 problems, we provide five numerical scores on a $[0, 1]$ scale:

\begin{itemize}  \item \textbf{Impact}: Mathematical significance and breadth of consequences;
  \item \textbf{Underexplored}: Degree to which the problem has received insufficient attention;
  \item \textbf{Toolability}: Amenability to computational or automated tools;
  \item \textbf{Formality}: Readiness for formal statement in a proof assistant;
  \item \textbf{AI Fit}: Overall suitability for current AI approaches.
\end{itemize}

\subsection{Domain and Subdomain Taxonomy}

OPA covers three top-level domains, each subdivided into fine-grained subdomains.
Mathematics is the largest domain (998 problems), encompassing subdomains from number theory and combinatorics to algebraic geometry and topology.
Theoretical computer science (541 problems) covers complexity theory, algorithms, cryptography, quantum computing, and learning theory.
Mathematical physics (536 problems) spans quantum field theory, statistical mechanics, general relativity, condensed matter theory, and quantum information.

% ============================================================
\section{Verification Framework}
\label{sec:verification}

\subsection{Multi-Layer Architecture}

OPA's verification framework operates at four layers:

\begin{description}  \item[L1 -- Schema Validation:] Every problem YAML file is validated against the JSON Schema (draft 2020-12) during CI. This ensures structural correctness: required fields are present, enumerations are valid, tier assignments are consistent with verification profiles.
  \item[L2 -- Verification Contracts:] Each problem has an associated \emph{verification contract} (a YAML file) that specifies: (a) the backend (\texttt{python\_checker} or \texttt{lean4}); (b) the task type (e.g., \texttt{statement\_verification}, \texttt{conjecture\_check\_range}, \texttt{proof\_check}, \texttt{counterexample\_search}); (c) resource limits (timeout, memory); and (d) success criteria in natural language.
  \item[L3 -- Executable Checkers:] For 101 problems, we provide Python checker scripts that implement specific computational verifications. Each checker exposes a \texttt{verify(params: dict) -> dict} interface, returning a structured result with \texttt{status} (\texttt{pass}/\texttt{fail}), \texttt{summary}, and \texttt{details}. Examples include range-checking the Collatz conjecture, verifying Goldbach representations up to a bound, and searching for counterexamples to graph-theoretic conjectures.
  \item[L4 -- Formal Verification:] For problems with Lean\,4 formalization stubs, the contract specifies \texttt{lean4} as the backend and \texttt{proof\_check} as the task type. A completed formal proof can be verified by the Lean type checker.
\end{description}

\subsection{Contract Specification}

A verification contract is a lightweight YAML document linking a problem to its verification infrastructure.
Listing~\ref{lst:contract} shows an example contract for the Collatz conjecture:

\begin{figure}[H]
\small
\begin{verbatim}
problem_id: opa.mathematics.number-theory.collatz-conjecture
backend: python_checker
task_type: conjecture_check_range
checker:
  file: verifiers/checkers/math/collatz_checker.py
  function: verify
resource_limits:
  timeout_seconds: 600
  memory_mb: 512
success_criteria: >
  Verify that the Collatz sequence reaches 1
  for all integers in [1, upper_bound]
parameters:
  upper_bound: 1000000
\end{verbatim}
\caption{Example verification contract for the Collatz conjecture.}
\label{lst:contract}
\end{figure}

\subsection{Checker Interface}

All Python checkers follow a uniform interface:

\begin{verbatim}
def verify(params: dict) -> dict:
    # Returns {"status": "pass"|"fail",
    #          "summary": str,
    #          "details": dict}
\end{verbatim}

This standardized interface enables automated batch execution.
The 101 checkers span multiple verification strategies: range checking of conjectures over integer intervals, counterexample search in combinatorial structures, bound verification for numerical conjectures, and witness construction for existential claims.

\subsection{Contract and Task Type Distribution}

Of the 2{,}075 contracts, 1{,}934 (93.2\%) use the \texttt{python\_checker} backend and 141 (6.8\%) use the \texttt{lean4} backend.
Task types are distributed as follows: statement verification (1{,}843), conjecture range checking (80), proof checking (141), counterexample search (5), bound verification (4), and witness search (2).

% ============================================================
\section{OPA-Bench}
\label{sec:bench}

\textsc{OPA-Bench} is a curated benchmark of 100 open problems designed for standardized evaluation of AI mathematical reasoning systems on unsolved questions.
Problems are selected from the full OPA dataset based on a composite \emph{bench score} that aggregates statement precision, solution checkability, machine actionability, and AI fitness.
The 100 highest-scoring problems are partitioned into three difficulty tiers by tertile:

\begin{itemize}  \item \textbf{Bronze} (33 problems): The most machine-tractable open problems. These have high statement precision, computational or proof-assistant checkability, and the highest AI fit scores. Examples: Collatz conjecture, Hadamard conjecture, Erd\H{o}s--Straus conjecture.
  \item \textbf{Silver} (33 problems): Problems with clear formal statements and partial machine checkability, but requiring more sophisticated approaches. Examples: Twin prime conjecture, Goldbach conjecture, Kakeya conjecture.
  \item \textbf{Gold} (34 problems): The most challenging benchmark problems, where machine actionability is lowest within the top 100. Examples: BPP vs.\ P, PPT bound entanglement, cosmic censorship conjecture.
\end{itemize}

Table~\ref{tab:comparison} compares OPA-Bench with existing mathematical reasoning benchmarks.

\begin{table}[H]
\centering
\caption{Comparison of OPA-Bench with existing mathematical reasoning benchmarks.}
\label{tab:comparison}
\small
\begin{tabular}{@{}lrccc@{}}
\toprule
\textbf{Benchmark} & \textbf{Problems} & \textbf{Type} & \textbf{Answers Known} & \textbf{Verification} \\
\midrule
MATH & 12{,}500 & Competition & Yes & Answer matching \\
GSM8K & 8{,}500 & Grade school & Yes & Answer matching \\
MiniF2F & 488 & Formal theorems & Yes & Lean/Isabelle \\
PutnamBench & 1{,}697 & Competition & Yes & Lean\,4 \\
\midrule
\textbf{OPA-Bench} & \textbf{100} & \textbf{Open problems} & \textbf{No} & \textbf{Multi-layer} \\
\bottomrule
\end{tabular}
\end{table}

The fundamental difference is that OPA-Bench does not have ground-truth answers.
Evaluation on OPA-Bench measures \emph{verifiable partial progress}---whether an AI system can produce outputs that pass the verification contract (e.g., extending a computational check to a larger range, finding a counterexample, or generating a Lean proof of a subresult).

\subsection{Domain Composition}

The 100 OPA-Bench problems span three domains: mathematics (50), theoretical computer science (30), and mathematical physics (20), reflecting a deliberate emphasis on problems where machine verification is most feasible.

% ============================================================
\section{Dataset Statistics}
\label{sec:stats}

\subsection{Domain Distribution}

Table~\ref{tab:domains} summarizes the distribution of problems across the three top-level domains.

\begin{table}[H]
\centering
\caption{Distribution of problems across domains.}
\label{tab:domains}
\begin{tabular}{@{}lrr@{}}
\toprule
\textbf{Domain} & \textbf{Count} & \textbf{Percentage} \\
\midrule
Mathematics & 998 & 48.1\% \\
Theoretical Computer Science & 541 & 26.1\% \\
Mathematical Physics & 536 & 25.8\% \\
\midrule
\textbf{Total} & \textbf{2{,}075} & \textbf{100\%} \\
\bottomrule
\end{tabular}
\end{table}

\subsection{Tier Distribution}

\begin{table}[H]
\centering
\caption{Distribution of problems across tiers.}
\label{tab:tiers}
\begin{tabular}{@{}lrrp{7.5cm}@{}}
\toprule
\textbf{Tier} & \textbf{Count} & \textbf{\%} & \textbf{Description} \\
\midrule
Tier 1 & 175 & 8.4\% & Core machine-actionable: high precision, checkable, actionable \\
Tier 2 & 106 & 5.1\% & Extended: precise statement, limited machine verification \\
Tier 3 & 1{,}794 & 86.5\% & Candidate: needs further structuring \\
\bottomrule
\end{tabular}
\end{table}

The dominance of Tier 3 problems reflects the inherent difficulty of unsolved mathematical problems: most require expert review for solution verification and are not yet amenable to fully automated approaches.
The 175 Tier 1 problems represent the immediate frontier for AI-driven mathematical research.

\subsection{Verification Profile Distribution}

Table~\ref{tab:verif-profile} shows the distribution along each axis of the verification profile.

\begin{table}[H]
\centering
\caption{Verification profile distribution across the three axes.}
\label{tab:verif-profile}
\begin{tabular}{@{}llrr@{}}
\toprule
\textbf{Axis} & \textbf{Value} & \textbf{Count} & \textbf{\%} \\
\midrule
Statement Precision & High   & 281   & 13.5\% \\
                    & Medium & 1{,}791 & 86.3\% \\
                    & Low    & 3     & 0.1\% \\
\midrule
Solution Checkability & Proof Assistant  & 144   & 6.9\% \\
                      & Computational    & 88    & 4.2\% \\
                      & Mixed            & 49    & 2.4\% \\
                      & Expert Review    & 1{,}794 & 86.5\% \\
\midrule
Machine Actionability & High   & 77    & 3.7\% \\
                      & Medium & 157   & 7.6\% \\
                      & Low    & 1{,}841 & 88.7\% \\
\bottomrule
\end{tabular}
\end{table}

\subsection{Score Distributions}

For the 312 problems with numerical scores, Table~\ref{tab:scores} reports summary statistics.

\begin{table}[H]
\centering
\caption{Score distributions for 312 annotated problems (scale 0--1).}
\label{tab:scores}
\begin{tabular}{@{}lcc@{}}
\toprule
\textbf{Score} & \textbf{Mean} & \textbf{Median} \\
\midrule
Impact        & 0.798 & 0.80 \\
Underexplored & 0.381 & 0.40 \\
Toolability   & 0.394 & 0.40 \\
Formality     & 0.803 & 0.85 \\
AI Fit        & 0.281 & 0.25 \\
\bottomrule
\end{tabular}
\end{table}

The high mean impact (0.80) reflects that the scored subset was drawn from well-known, significant problems.
The low mean AI fit (0.28) quantitatively captures the difficulty of unsolved problems for current AI methods.
The gap between formality (0.80) and AI fit (0.28) suggests that many problems are sufficiently well-formalized but remain beyond current AI capabilities---precisely the gap OPA is designed to help close.

\subsection{Top Subdomains}

Table~\ref{tab:subdomains} lists the 15 most represented subdomains.

\begin{table}[H]
\centering
\caption{Top 15 subdomains by problem count.}
\label{tab:subdomains}
\small
\begin{tabular}{@{}lr@{\hspace{2em}}lr@{}}
\toprule
\textbf{Subdomain} & \textbf{Count} & \textbf{Subdomain} & \textbf{Count} \\
\midrule
Number Theory          & 203 & Topology              & 70 \\
Algorithms             & 124 & General Relativity    & 61 \\
Quantum Information    & 116 & Cryptography          & 60 \\
Complexity Theory      & 107 & Probability           & 59 \\
Combinatorics          & 97  & Quantum Computing     & 55 \\
Analysis               & 93  & String Theory         & 53 \\
Algebra                & 92  & Graph Theory          & 51 \\
Quantum Field Theory   & 87  & & \\
\bottomrule
\end{tabular}
\end{table}

\subsection{Problem Status}

Of the 2{,}075 problems, 2{,}065 (99.5\%) are currently open, 4 are solved, 3 are partially solved, and 3 are disproved.
The small number of non-open problems serves as ground truth for validating the status-tracking pipeline: these problems were open when first entered and their status was later updated based on new results.

\subsection{OPA-Bench Difficulty Breakdown}

\begin{table}[H]
\centering
\caption{OPA-Bench difficulty distribution.}
\label{tab:bench-difficulty}
\begin{tabular}{@{}lrrrrr@{}}
\toprule
\textbf{Difficulty} & \textbf{Count} & \textbf{Math} & \textbf{TCS} & \textbf{Math-Phys} \\
\midrule
Bronze & 33 & 32 & 0 & 1 \\
Silver & 33 & 18 & 12 & 3 \\
Gold   & 34 & 0  & 18 & 16 \\
\midrule
\textbf{Total} & \textbf{100} & \textbf{50} & \textbf{30} & \textbf{20} \\
\bottomrule
\end{tabular}
\end{table}

The domain composition shifts across difficulty levels: Bronze is dominated by mathematics (particularly number theory and combinatorics), where computational verification is most natural.
Gold is dominated by theoretical computer science and mathematical physics, where machine actionability is more limited.

% ============================================================
\section{Lean\,4 Formalization}
\label{sec:lean4}

OPA includes 312 Lean\,4 formalization stubs, covering 15.0\% of the dataset.
Each stub is a minimal Lean\,4 file that:
\begin{enumerate}  \item Documents the problem with its OPA identifier, domain, status, and informal statement in a block comment;
  \item Declares the problem as an \texttt{axiom} with \texttt{True} as the target type, serving as a placeholder for the actual formal statement and proof;
  \item Provides a URL back to the OPA web entry.
\end{enumerate}

While these stubs are intentionally minimal---they do not formalize the mathematical content of the statements---they serve several purposes.
First, they establish the scaffolding for a formalization effort: the file structure, naming convention, and metadata are in place, reducing the friction for contributors who wish to add formal statements.
Second, 141 problems have verification contracts with \texttt{lean4} as the backend and \texttt{proof\_check} as the task type, creating a pipeline where a formalized proof can be automatically verified.
Third, the stubs identify problems where formalization is explicitly \emph{wanted} (via the \texttt{formalization.wanted: true} flag in the problem YAML), enabling prioritization of formalization efforts.

The formalization stubs complement recent work on autoformalization~\citep{wu2022autoformalization,glazer2024formalconjectures}, which could be applied to automatically upgrade the stubs from placeholder axioms to genuine Lean\,4 statements.

% ============================================================
\section{Limitations and Future Work}
\label{sec:limitations}

\paragraph{Coverage and Bias.}
The dataset's domain distribution reflects the availability of well-documented open problems in existing sources.
Areas with fewer accessible problem lists (e.g., applied mathematics, mathematical biology) are underrepresented.
The subdomain taxonomy is not fully standardized and may contain overlapping or inconsistent categories.

\paragraph{Status Currency.}
Mathematical problems can be solved at any time, and our status-tracking pipeline relies on periodic manual review.
While the schema includes \texttt{stale\_after} dates and review states to flag potentially outdated entries, there is inherent latency between a problem being solved and the dataset being updated.

\paragraph{Verification Depth.}
Only 101 of 2{,}075 problems have executable checkers, and the Lean\,4 stubs are placeholders.
Expanding the checker coverage and upgrading the formalization stubs to genuine formal statements are key priorities.

\paragraph{Score Subjectivity.}
The five-dimensional scores (impact, underexplored, toolability, formality, AI fit) are currently assigned through a semi-automated process and have been computed for only 312 problems.
Calibrating these scores through expert review and extending coverage to all problems remains future work.

\paragraph{Evaluation Methodology.}
OPA-Bench evaluates verifiable partial progress, but defining and measuring ``partial progress'' on open problems is inherently challenging.
We plan to develop more nuanced evaluation protocols, including multi-turn interaction frameworks where AI systems can query checkers iteratively.

\paragraph{Future Directions.}
We plan to: (1)~expand the dataset to 5{,}000+ problems with broader domain coverage; (2)~upgrade Lean\,4 stubs to full formal statements using autoformalization; (3)~develop a multi-turn verification protocol for evaluating AI research assistants; (4)~integrate with HuggingFace for periodic dataset snapshots; and (5)~build a community curation platform for expert review and annotation.

% ============================================================
\section{Conclusion}
\label{sec:conclusion}

We have presented \textsc{OpenProblemAtlas}, a dataset of 2{,}075 unsolved mathematical problems with structured annotations, verification contracts, executable checkers, and formalization stubs.
Together with the \textsc{OPA-Bench} benchmark, OPA provides the first infrastructure for evaluating AI systems on genuine open questions in mathematics, theoretical computer science, and mathematical physics.
By shifting evaluation from solved problems to unsolved ones, OPA aims to redirect AI mathematical reasoning research toward the ultimate goal: contributing to human mathematical knowledge.

% ============================================================
\bibliographystyle{plainnat}
\bibliography{references}

\end{document}
